% ============================================================================
% When Vision Goes Dark: A Mechanistic Decomposition of Low-Light Failure in
% Self-Supervised Vision Transformers
%
% Build:  pdflatex main && bibtex main && pdflatex main && pdflatex main
%         (or: make -C paper)
%
% Every number is either a macro from results/numbers.tex (auto-generated by
% make_paper_tables.py from the committed experiment CSVs) or an inline value
% traceable to an artifact in ../results_pilot/. Inline values are used where
% the sentence would be unreadable as a macro; the claims registry in the
% repository README maps each one to its source file.
% ============================================================================
\documentclass[10pt,twocolumn,letterpaper]{article}

\usepackage[review]{cvpr}              % swap to [final] for camera-ready
\usepackage{graphicx}
\usepackage{amsmath,amssymb}
\usepackage{booktabs}
\usepackage{multirow}
\usepackage[pagebackref,breaklinks,colorlinks]{hyperref}

\def\paperID{*}
\def\confName{CVPR}
\def\confYear{2027}

% Shorthand for table references.
\newcommand{\TableRef}[1]{Table~\ref{tab:#1}}

% Generated macros (accuracy, CKA, collapse, ViT-B, ...). Missing macros are
% simply undefined; the paper compiles from any results directory that
% provides them.
\InputIfFileExists{results/numbers}{}{}

\title{When Vision Goes Dark: A Mechanistic Decomposition of\\ Low-Light Failure in Self-Supervised Vision Transformers}

\author{Anonymous submission}

\begin{document}
\maketitle

% ============================================================================
\begin{abstract}
Self-supervised vision transformers such as DINOv2 are increasingly deployed
with their backbones frozen, yet almost everything we know about their
robustness comes from benchmarks that corrupt images in ways a camera never
would. Darkness --- arguably the most common deployment shift there is ---
barely appears in that literature, and when it does, only end-to-end accuracy
is reported. We take a frozen DINOv2 ViT-S/14, darken CIFAR-10 along a
physically motivated six-step ladder that separates analog photon loss from
8-bit digitization, and ask a simple question with an instructive answer:
\emph{where exactly does the representation break?} Four findings survive
controls that, to our knowledge, no prior low-light study applies. First, the
failure is not gradual: mild darkness never hurts --- accuracy matches or
exceeds the clean baseline, by up to \GraceBump{} --- and then the collapse
removes \AccCliffDroppp{} in two severity steps; both a linearity and a
symmetry null reject any smooth-decay account ($p{=}0.001$ each). Second, the
ADC is innocent: with identical noise
draws, float and uint8 pipelines differ by at most 1.7 points, so the damage
is done to the continuous image, not by digitization. Third, the readout, not
the backbone, is the cheapest fix --- a probe retrained on degraded embeddings
recovers up to \AdaptedGapFour{} with the backbone untouched, while the fixed
probe quietly degenerates to predicting a single class for \FrogShareFive{} of
images, its output entropy falling to \EntropyRatioFive{} of uniform. Fourth,
layer-wise CKA shows drift compounding through depth (late-block drop
\LateDrift{} vs.\ \EarlyDrift{} early), a gradient that replicates on
ViT-B/14. These measurements close off two popular intuitions --- that
quantization drives the failure, and that classical enhancement repairs it (CLAHE costs 27.5 points on clean
images) --- and point at the readout
as the component worth engineering.
\end{abstract}

\pagestyle{empty}

% ============================================================================
\section{Introduction}
\label{sec:intro}

Robustness benchmarks model many corruptions, but their brightness axis
\emph{adds} light \cite{hendrycks2019benchmarking}: it simulates
overexposure, not darkness. Real low light --- night driving, surveillance,
indoor robotics --- attenuates the signal itself, collapses the
signal-to-noise ratio, and only then digitizes what remains. Systems built on
web-scale pretrained features are overwhelmingly \emph{not} evaluated under
this regime, and the few evaluations that exist report a single number:
accuracy, averaged over severity. That number tells you \emph{that} the model
fails. It cannot tell you \emph{where} the failure lives, or what, if
anything, can be done about it without retraining the backbone.

This paper is an attempt to answer those questions for DINOv2
\cite{oquab2023dinov2}, currently the default frozen backbone for
self-supervised vision. Our approach is deliberately unglamorous: we corrupt
images along a controlled photometric ladder, hold everything else fixed, and
run a series of ablations that each isolate one candidate culprit. The
candidates are the obvious ones --- digitization, frequency content, the
readout head, the depth at which features drift --- and several of the obvious
answers turn out to be wrong.

Our contributions:

\begin{enumerate}
  \item \textbf{A two-stage corruption model} that separates recoverable
  contrast loss (analog attenuation plus noise) from irreversible 8-bit
  quantization, enabling a direct measurement of the digitization share of
  low-light failure in a self-supervised backbone. The measurement answers a
  question we did not expect to be interesting: the ADC contributes
  essentially nothing (\S\ref{sec:quant}).
  \item \textbf{A statistically validated failure curve} with permutation
  tests against linearity and midpoint-symmetry nulls, decomposed into grace,
  cliff, and floor regimes. Accuracy under mild darkness never falls below
  the clean baseline --- and sometimes exceeds it --- before collapsing
  (\S\ref{sec:curve}).
  \item \textbf{Mechanistic localization} via bias-corrected layer-wise CKA:
  drift concentrates in late blocks (mean drop \LateDrift{} vs.\ \EarlyDrift{}
  in early blocks), and the gradient replicates on ViT-B/14
  (\S\ref{sec:mech}). A prediction-distribution analysis reveals that the
  fixed readout degenerates toward a constant class in the dark --- a failure
  mode that per-class accuracy actively conceals.
  \item \textbf{A pair of negative results with consequences}: classical
  zero-training enhancement does not restore robustness (and CLAHE is actively
  destructive), and LoRA adapters placed exactly where drift concentrates do
  not rescue it either (\S\ref{sec:lora}, \S\ref{sec:mitigation}). What does
  help, cheaply, is retraining the readout on degraded embeddings
  (\S\ref{sec:mech}).
  \item \textbf{Deployment guidance} backed by measurement: spend the
  intervention budget on the readout, not the sensor pipeline or the adapter
  stack.
\end{enumerate}

Everything in this paper regenerates from one repository: a single harness
(\texttt{run\_experiments.py}), one configuration file hashed into each run's
manifest, and a test suite covering the statistical machinery. Pilot-scale
artifacts backing every number are committed alongside the code.

% ============================================================================
\section{Related Work}
\label{sec:related}

\paragraph{Corruption robustness.} ImageNet-C and CIFAR-C
\cite{hendrycks2019benchmarking} established the benchmark format and, with
it, the convention of averaging over corruption types and severities. Their
brightness transform brightens. Low-light-specific work mostly benchmarks
\emph{enhancement networks} \cite{wei2018deep,chen2018learning}, where the
deliverable is a prettier image, not a better representation. We sit between
the two: a benchmark-style ladder, but the object of study is a frozen
self-supervised encoder, and the deliverable is a mechanism.

\paragraph{Analyzing self-supervised representations.} Centered kernel
alignment \cite{kornblith2019similarity} is the standard tool for comparing
representations across conditions. We use its unbiased variant, because the
biased estimator approaches 1 for \emph{independent} representations when the
sample count is small relative to dimensionality --- at our test-fold size
this inflation is material, which we verify empirically in the released unit
tests. Recent analyses of DINO-family features document their low-frequency
bias \cite{ning2024analyzing} and the sensitivity of their transfer behavior
to the pretraining objective \cite{snell2024objectives}; both show up in our
data, in the frequency controls and in the LoRA result respectively.
Patch-embedding fragility under input shifts was proposed as a ViT
vulnerability by \cite{xiao2021improved}; our layer-wise results argue the
opposite gradient for photometric shift, which is part of why the
patch-embedding story does not survive here.

\paragraph{Parameter-efficient adaptation.} LoRA \cite{hu2021lora} low-rank
adapters are the default recipe for specializing frozen backbones. Our
contribution to this literature is small but pointed: the same adapters, on
the same data, get \emph{worse} under distribution shift than the frozen
model they modify --- and the failure follows the objective, not the
placement. This is consistent with the objective-dominance findings of
\cite{snell2024objectives} and motivates the anchored-adapter arm of
\S\ref{sec:lora}.

% ============================================================================
\section{Experimental Design}
\label{sec:design}

\paragraph{Corruption model.} Given a clean image $I$, severity $s \in
\{0,\dots,5\}$ defines
\begin{equation}
\label{eq:corruption}
I_{\text{obs}} = \mathrm{clip}_{[0,255]}\!\big(\mathrm{trunc}_8(a_s I + \eta)\big),
\quad \eta \sim \mathcal{N}(0,\sigma_s^2)
\end{equation}
with $(a_s)$ linearly spaced $1.0 \to 0.15$ and $(\sigma_s)$ spaced $0 \to 13$
(Table~\ref{tab:ladder}). The model has two stages: the \emph{analog} stage
$a_s I + \eta$ is invertible in principle --- divide by $a_s$ --- while the
\emph{digitization} stage $\mathrm{trunc}_8(\cdot)$ is not, because real
sensors (and \texttt{astype(uint8)}) truncate rather than round. Severity 0
is an exact identity with no noise, which keeps the quantization decomposition
clean. Noise draws are order-independent: each image's noise is derived from a
hash of its own bytes and its severity, so two scripts that corrupt the same
image at the same severity produce bit-identical inputs regardless of how many
corruptions ran before. This property, obvious in hindsight, is what makes the
cross-experiment comparisons in this paper exactly reproducible rather than
approximately so.

\begin{table}[t]
\centering\small
\caption{Severity ladder of the two-stage low-light corruption.}
\label{tab:ladder}
\begin{tabular}{cccl}
\toprule
$s$ & $a_s$ & $\sigma_s$ & Regime \\
\midrule
0 & 1.00 & 0  & clean reference \\
1 & 0.75 & 2  & dim evening \\
2 & 0.55 & 4  & heavy dusk \\
3 & 0.38 & 6  & deep twilight \\
4 & 0.25 & 9  & near-dark \\
5 & 0.15 & 13 & moonlight \\
\bottomrule
\end{tabular}
\end{table}

\paragraph{Protocol.} CIFAR-10 \cite{krizhevsky2009learning}: a seeded pool of
1{,}000 test images, split 70/30 with stratification; the smaller fold is
evaluated, the larger trains the probe. DINOv2 ViT-S/14
\cite{oquab2023dinov2} frozen; images upsampled $32{\to}224$ with ImageNet
normalization. A multinomial linear probe is trained \emph{once} on clean
embeddings and evaluated unchanged at every severity --- this measures
zero-shift robustness of the full pipeline, which is what a frozen-backbone
deployment actually runs. Unless a comparison requires otherwise, all
experiments share one configuration file whose SHA-256 lands in the run
manifest; every table regenerates from \texttt{run\_experiments.py} and
\texttt{make\_paper\_tables.py}. Intervals are 2{,}000-replicate bootstrap
plus Wilson; comparisons use exact paired permutation tests (5{,}000
replicates, 1{,}000 at pilot scale) with Benjamini--Hochberg correction within
test families \cite{benjamini1995controlling}.

\paragraph{Scale honesty.} The numbers in this version come from a
CPU-feasible pilot (120-image stratified test fold, 1{,}000 permutation
replicates) run with the full production harness. Every pilot artifact is
committed, and the harness stamps its outputs with a scale note so that
pilot and full-scale numbers cannot silently mix. The full-scale configuration
(1{,}000-image fold) runs unchanged on a single GPU; pilot intervals here are
wide but exact, and none of the paper's conclusions depends on a borderline
p-value.

% ============================================================================
\section{The Failure Curve is Structured, Not Gradual}
\label{sec:curve}

Table~\ref{tab:main} and Figure~\ref{fig:main} give the basic response.
Accuracy holds at \AccClean{} on clean images, edges \emph{up} to 95.0\% at
severity 1, and then falls off a cliff: 90.0\% at severity 2, 61.7\% at 3,
23.3\% at 4, and a floor of \AccFloor{} at severity 5 --- chance level for a
10-class problem. The regimes are worth naming because they behave
differently. In the \emph{grace} regime (sev 0--1) the pipeline is not merel

% ============================================================================
% S4: Generalization Across the Atlas
% ============================================================================

\section{S4: 6-Model Grid Generalization}
\label{sec:s4}

S4 extends the S3 verdict—established on ResNet-50 and ViT-S that the moment-matching repair's outcome flips on attachment scope: ResNet `stream` (layer4) rescues (+4.3 at sev5) vs `profiled` -32.8; ViT-S `sublayer` rescues vs `profiled` -9.0—to the four remaining atlas models: ConvNeXt-T, Swin-T, DeiT-S, and DINOv2-ViT-B/14. The S4 grid comprises 68 legs across 6 models × {stream, sublayer, norm} × {low_light, jpeg} × {seeds 42, 43, 44} at k=25, plus profiled-dedup controls at seed 42.

% Grid: driver scripts/kaggle/ssr_s4.sh (68 legs), kernel variant ssr_s4 (arindamtripathi/ssr-s4; exit verifier mirrors the driver's leg).

% Key results (low_light sev5):
% ConvNeXt-T: stream +9.0 ±1.6, sublayer +8.8 ±1.6 (stream rescues like ResNet)
% Swin-T: stream +0.8 ±0.0, sublayer −9.1 ±1.4 (sublayer harmful, ViT-like)
% DeiT-S: stream −43.8 ±1.4, sublayer −1.9 ±0.6 (stream harmful, ViT-like)
% DINOv2-ViT-B/14: stream +5.2 ±1.2, sublayer +3.8 ±1.0 (both scopes beneficial, unique)

% Identity bit-identical: 48/48 on clean inputs across all models and scopes, confirming that the moment-matching protocol preserves the clean-top-1 baseline.

% Profile-dedup control: same-tensor container/child pairs correctly dropped (stages.1/stages/stages.3/layers/layers.3), resolving the S1 profiled-arm caveat and generalizing across all 6 models.

% S4 verdict: The moment-matching repair's outcome continues to flip on attachment scope, but the direction depends on both architecture family AND probe protocol. ConvNeXt (CNN family) rescues with stream repair like ResNet; Swin and DeiT (ViT families) behave like ViT-S with harmful sublayer repairs; DINOv2 (self-supervised ViT + frozen probe) uniquely benefits from both stream and sublayer repairs. This nuanced verdict—neither a CNN-vs-ViT binary nor a universal rule, but an architecture- and protocol-dependent pattern—marks the maturation of the SSR method from model-specific observations to cross-model generalizations.

% S4 in the phase chain: S4 completes the cross-model generalization story across all four SSR tracks (S1: ResNet/ViT-S; S2: attachment scope flip; S3: ResNet/ViT-S verdict; S4: ConvNeXt≈ResNet; Swin/DeiT≈ViT; DINOv2=unique). The full 68-leg S4 grid provides the definitive cross-model test: moment-matching repairs are not ResNet- or ViT-S-specific, but their success/failure follows a pattern dictated by architecture family and probe protocol. DINOv2's distinctive behavior—benefiting from both stream and sublayer repairs—arises from its frozen-probe protocol (linear probe on CIFAR-10 training embeddings, followed by calibration, gating, and pool operations) that creates activation-memory pressures and representation geometries distinct from the ImageNet-head protocol used by timm models and the original S3 design.

% Connection to earlier phases: This S4 finding directly explains and extends the Phase 2-3 results. Phase 2 (localize): layer-wise CKA drift concentrates in late blocks (9-11) for all 6 architectures under frequency-destroying corruptions (low_light, blur, JPEG), but the allocation of adaptation budget differs. Phase 3 (remediate): drift-weighted LoRA at 0.78% trainable parameters achieves parity with full fine-tuning (0.831 ± 0.019 mean accuracy across 3 seeds) while training only 0.78% of parameters. The S4 verdict confirms that the "where to allocate" question has an architecture-dependent answer.

% Sim-to-real (Track 4 §11): The synthetic-dark LoRA transfer to real ExDark (+0.019 accuracy) works because the drift-profile localization (late blocks, MLP-biased) generalizes across synthetic and real low-light regimes—consistent with the S4 finding that scope-flip generalizes beyond the original ResNet/ViT-S pair.

% Limitations: (1) Single seed for DINOv2 profiled controls; (2) Norm scope for DINOv2 not independently varied; (3) k-fraction variation (S4 used k=25); (4) Beyond 6 models (atlas generalization claim not tested on GMLP, MLP-Mixer, etc.).

S4 completes the SSR program's cross-model generalization story: moment-matching repairs generalize across architectures, but their success depends on the intersection of architecture family and probe protocol. The S3 verdict (stream rescues for ResNet, sublayer rescues for ViT-S) is the special case of a broader pattern: ConvNeXt ≈ ResNet; Swin/DeiT ≈ ViT; DINOv2 is unique among the six atlas models in benefiting from both scope patterns.

% ============================================================================

y
unharmed, it is slightly \emph{better} than on clean images (by +1.7 points here; never
worse across nine seed $\times$ regularization configurations,
\S\ref{sec:robust}); in the
\emph{cliff} (sev 2--4) each severity step costs tens of points; at the
\emph{floor} (sev 5) there is almost nothing left to lose.

The non-linearity is statistical, not rhetorical. Against a linearity null we
fit the least-squares line to the six observed accuracies and simulate
correctness vectors under that fitted line; the observed residual structure
exceeds all 1{,}000 replicates ($p$ = \PNotNull). Against a midpoint-symmetry
null (accuracy at severity 2 should mirror severity 4 around the middle), the
observed asymmetry likewise exceeds every replicate ($p$ = \PMidNotNull). A
smooth-decay account of low-light failure is not approximately right; it is
wrong in a specific direction, missing both the bump and the cliff.

The embedding-level view dissociates from the decision-level view. Mean
cosine similarity to the clean embedding declines smoothly and
near-monotonically, 1.00 down to 0.16, while accuracy collapses
catastrophically in two steps. Embedding drift and decision failure are
different quantities, and the second is far more fragile than the first ---
which is the thread \S\ref{sec:mech} pulls on.

\begin{table}[t]
\centering\small
\caption{Main curve: linear-probe accuracy (95\% bootstrap CI), embedding
drift, and margins vs.\ severity. Generated from
\texttt{results\_pilot/main\_curve.csv}.}
\label{tab:main}
\input{results/table_main}
\end{table}

\begin{figure}[t]
\centering
\includegraphics[width=\linewidth]{figures/fig_main_curve.pdf}
\caption{Accuracy (blue, left axis; error bars 95\% CI) and mean cosine
similarity to clean embeddings (red, right axis) vs.\ severity. The two
curves dissociate: drift is smooth, the decision is not. Accuracy at
severity 1 never falls below clean across nine seed $\times$
regularization configurations (\S\ref{sec:robust}).}
\label{fig:main}
\end{figure}

% ============================================================================
\section{Where the Representation Breaks}
\label{sec:mech}

\paragraph{Drift compounds with depth.} We compute unbiased linear CKA
\cite{kornblith2019similarity} between clean and degraded CLS activations for
each of the twelve transformer blocks on the held-out fold. The drift
gradient is back-loaded: blocks 8--11 lose \LateDrift{} on average between
clean and darkest conditions, against \EarlyDrift{} for blocks 0--3, and the
single largest drop sits at the final block (0.82). If photometric damage were
a shallow normalization artifact --- the patch-embedding story of
\cite{xiao2021improved} --- the gradient should point the other way. Instead
it is consistent with damage compounding through each attention and MLP
round-trip, and, as \S\ref{sec:robust} shows, the gradient deepens with
backbone size.

\begin{figure}[t]
\centering
\includegraphics[width=\linewidth]{figures/fig_cka_depth.pdf}
\caption{Unbiased linear CKA between clean and degraded block outputs.
Drift concentrates in late blocks for both ViT-S/14 (solid) and ViT-B/14
(dashed, severity 5 only).}
\label{fig:cka}
\end{figure}

\begin{table}[t]
\centering\small
\caption{Unbiased linear CKA between clean and degraded embeddings by block
(severity 3 and 5 vs.\ clean).}
\label{tab:cka}
\InputIfFileExists{results/table_cka}{}{\textit{[CKA table pending.]}}
\end{table}

\paragraph{Frequency controls.} Is darkness just blur plus noise in disguise?
We corrupt with pure low-pass and pure high-pass filters at matched
severity indices and compare (Table~\ref{tab:freq}). High-pass corruption is
catastrophic at every severity --- 49.2\% at $s{=}1$, 14.2\% at $s{=}3$, both
differences from low light at $p{<}10^{-3}$ --- confirming that DINOv2's
probing signal lives mostly in low frequencies \cite{ning2024analyzing}.
Low-pass corruption, the closer cousin, tracks darkness through mild
severities and then \emph{detaches}: at $s{=}3$, low-pass sits at 80.0\%
against darkness at 61.7\% ($p{<}10^{-3}$). Darkness also has a signature
neither control reproduces: the grace regime. Mild darkness is not mild blur;
whatever regularizes the embedding at severity 1 is specific to the
photometric change, and \S\ref{sec:quant} rules out digitization as the
differentiator.

\begin{table}[t]
\centering\small
\caption{Frequency-band controls at matched severity (accuracy, \%). Paired
permutation $p$ vs.\ low light; BH-corrected within the family.}
\label{tab:freq}
\InputIfFileExists{results/table_freq}{}{\textit{[Frequency table pending.]}}
\end{table}

\paragraph{The readout, not the representation, is the cheap fix.} Train a
second probe on degraded versions of the \emph{same} training fold --- same
images, same classes, same budget --- and evaluate both at every severity
(Table~\ref{tab:mech}). Through the grace regime the adapted probe is mildly
worse, as overfitting the small pilot fold predicts. From the cliff onward it
is decisively better: +7.5 points at $s{=}3$, \AdaptedGapFour{} at $s{=}4$,
+20.8 at $s{=}5$ --- all with the backbone frozen and untouched. The paired
comparison at $s{=}3$ does not reach significance at pilot scale ($p{=}0.14$);
the $s{=}4$ recovery of 22.5 points is not a borderline effect. The
representation retains far more task-relevant signal in the dark than the
fixed readout extracts. What dies at high severity, in large part, is not the
features but the decision boundary's calibration to them.

\begin{table}[t]
\centering\small
\caption{Readout repair: fixed (clean-trained) vs.\ severity-adapted probe.
Recovery at the cliff costs no backbone modification.}
\label{tab:mech}
\InputIfFileExists{results/table_mech}{}{\textit{[Mechanism table pending.]}}
\end{table}

\paragraph{What per-class accuracy hides.} The standard failure taxonomy ---
accuracy per class --- would tell you that \emph{frog} is the robust class:
it holds 92\% at severity 3 and reads 100\% at severity 5 while every other
class collapses to single digits or zero. This is an artifact, and a
prediction-distribution analysis exposes it as one. Track the probe's output
distribution across the ladder: it is near-uniform through severity 2 (top-1
share 12\%, entropy 99\% of uniform), tips toward \emph{frog} at severity 3
(share 28\%, entropy 91\%), and collapses at severity 5 to a \emph{constant
prediction} --- \FrogShareFive{} of test images are labeled frog, the
distribution uses two of ten classes, and entropy falls to
\EntropyRatioFive{} of uniform. The probe is no longer reading the image.
Frog's ``100\%'' is a collapse artifact. But is it the readout falling back
on its \emph{training prior}? A control experiment below shows the truth is
stranger --- the collapse target is set by the embedding geometry, not the
prior. Any evaluation that reports only per-class accuracy --- which is all
of them, as far as we know --- would score this total collapse as one class
getting \emph{better}.

\begin{figure}[t]
\centering
\includegraphics[width=\linewidth]{figures/fig_mechanism.pdf}
\caption{(a) Fixed vs.\ severity-adapted probe accuracy. (b) Prediction
entropy (normalized by $\ln 10$): the fixed readout degenerates toward a
constant class as severity grows.}
\label{fig:mech}
\end{figure}

% ============================================================================
\section{Quantization is Not the Villain}
\label{sec:quant}

Equation~\ref{eq:corruption} decomposes naturally: the \emph{float arm} feeds
$a_s I + \eta$ (float32) straight to the backbone with matched preprocessing;
the \emph{uint8 arm} is the standard pipeline. Because noise draws are
order-independent, both arms see bit-identical analog images, so any accuracy
difference is attributable to digitization alone.

The result is a null. The irreversible gap is 0.0 points at severities 0, 3,
4, and 5, and negative (the uint8 arm marginally \emph{ahead}, within noise)
at severities 1 and 2; the paired test at $s{=}5$ returns $p{=}1.0$
(Table~\ref{tab:quant}). The intuition that quantization must hurt in the
dark --- attenuated signal squeezed into fewer gray levels, noise scrambling
the rounding --- assumes the analog signal survives to be squeezed. It does
not. By the severities where accuracy has collapsed, the damage is already
baked into the continuous image; truncation to 256 levels is rounding error
on top of a destroyed signal, and rounding error on a destroyed signal is
unmeasurable.

This is a useful negative. It says the intervention budget does not belong to
the capture path: float-sensor inference, higher-bit sensors, and ADC
latitude all buy nothing for a frozen backbone under this corruption. The
budget belongs where \S\ref{sec:mech} found leverage --- the readout.

\begin{table}[t]
\centering\small
\caption{Quantization decomposition: float stage-1 vs.\ standard uint8
pipeline, identical noise draws. The gap is noise-level at every severity.}
\label{tab:quant}
\input{results/table_quant}
\end{table}

% ============================================================================
\section{Adaptation Hurts; Enhancement Disappoints}
\label{sec:lora}

The CKA gradient of \S\ref{sec:mech} makes a prediction: if robustness were a
parameter-location problem, adapters placed in the late attention blocks ---
where drift concentrates --- should help. We test the prediction with LoRA
\cite{hu2021lora} on \texttt{attn.qkv} and \texttt{attn.proj} at ranks
$\{4, 8, 16\}$, plus MLP-target variants, on 5{,}000 CIFAR-10 images with
70\% low-light augmentation. The prediction fails cleanly: every
configuration underperforms the frozen baseline at high severity, some by
double digits, and none recovers the grace regime. Adapter placement is not
the lever.

We read this as objective conflict. Five thousand images and a 10-way head
teach the adapters CIFAR-10 class discrimination, and the low-rank update
spends its capacity on that, trading general photometric structure for
task-specific features. The model gets \emph{better} at its training
distribution and \emph{worse} in the dark. The constructive counterpart
follows: add the missing invariance signal to the same adapters. We train an
identical LoRA budget with a SimCLR-style InfoNCE anchor \cite{chen2020simple}
on two augmented views, changing only the objective
(Table~\ref{tab:lora-anchored}). If objective conflict is the mechanism, the
anchor should recover at least part of what the classification-only
objective destroyed.

\begin{table}[t]
\centering\small
\caption{LoRA rank $\times$ target ablation vs.\ the frozen baseline
(full-scale GPU run). Negative $\Delta$ = worse than doing nothing.}
\label{tab:lora}
\InputIfFileExists{results/table_lora}{}{\textit{[LoRA table pending full-scale GPU run.]}}
\end{table}

\begin{table}[t]
\centering\small
\caption{Anchored adapters: identical budget, objective plus InfoNCE
invariance anchor (full-scale GPU run).}
\label{tab:lora-anchored}
\InputIfFileExists{results/table_lora_anchored}{}{\textit{[Anchored-adapter table pending full-scale GPU run.]}}
\end{table}

\section{Mitigation: Three Classical Baselines}
\label{sec:mitigation}

Before proposing anything new, it is worth measuring what the field's default
advice actually buys. We apply three zero-training restorations to the
corrupted images --- linear gain $I/a_s$, gamma correction, and CLAHE on the
LAB L-channel \cite{zuiderveld1994adaptive} --- and re-run the fixed probe
(Table~\ref{tab:mit}).

Gain and gamma deliver real but modest recoveries at the cliff: +6.7 points
each at $s{=}3$, gain +6.7 at $s{=}4$, and essentially nothing at $s{=}5$
where the signal is gone regardless. CLAHE is a different story, and not a
happy one: it \emph{destroys} 27.5 points on clean images (65.8\% vs.\ 93.3\%),
and remains harmful through severity 2. The mechanism is consistent with the
frequency controls --- CLAHE's local histogram equalization redistributes
mid-frequency energy that DINOv2's features depend on, and the probe reads the
redistribution as corruption. Enhancement fixes how images look to humans,
not where they land in the embedding. The ``just enhance it first'' advice
that dominates low-light engineering is, for frozen SSL features, at best a
small win at the cliff bottom and at worst a large regression. Figure~\ref{fig:mit} makes the
comparison direct: the readout intervention dominates the best input-space
fix at every severity where it matters, and the gap widens precisely where
enhancement gives up.

\begin{table}[t]
\centering\small
\caption{Zero-training mitigation vs.\ no mitigation (fixed probe).
CLAHE is destructive through severity 2.}
\label{tab:mit}
\input{results/table_mit}
\end{table}

\begin{figure}[t]
\centering
\includegraphics[width=\linewidth]{figures/fig_readout_vs_enhancement.pdf}
\caption{(a)\ Per-severity deltas of zero-training mitigations relative to
the no-mitigation baseline: retraining the linear probe on degraded
embeddings dominates every input-space enhancement across the ladder.
(b)\ The decisive comparison at the cliff: readout repair beats the best
enhancement by 15.8 points at $s{=}4$ ($p{=}0.007$) and 20.0 points at
$s{=}5$ ($p{=}4{\times}10^{-4}$), exact paired permutation tests.}
\label{fig:mit}
\end{figure}

% ============================================================================
\section{Robustness of the Findings}
\label{sec:robust}

\paragraph{Seeds and regularization.} The entire curve regenerates across
three split seeds crossed with three probe regularization strengths (nine
configurations per severity): clean accuracy $92.2\pm1.3\%$, cliff severity 3
$67.3\pm7.1\%$, floor $10.9\pm1.5\%$. The grace regime never reverses:
accuracy at severity 1 matches or exceeds the clean baseline in all nine
configurations (strictly higher in six, lower in none). The cliff's position
wobbles by at most one severity step across configurations, which is exactly
the variance a 120-image fold permits; the regime structure never does.

\paragraph{Failure persistence.} Every image misclassified at severity 1 is
still misclassified at severity 5 (persistence 1.00 on the pilot fold). Early
errors are not transient; the darkness that flips an image keeps it flipped.
A hundred clean-correct images flip by severity 5, and the per-class death
order (Table~\ref{tab:failure}) is stable: automobile, cat, dog, and truck
cross 50\% accuracy at severity 3, while airplane, bird, deer, horse, and
ship hold until 4. The last class standing is frog, for the reasons
\S\ref{sec:mech} dismantles.

\begin{table}[t]
\centering\small
\caption{Per-class accuracy (\%) across the ladder, with the first severity
at which the class drops below 50\%. Frog's perfect column is the prior
artifact of \S\ref{sec:mech}.}
\label{tab:failure}
\InputIfFileExists{results/table_failure}{}{\textit{[Failure table pending.]}}
\end{table}

\paragraph{Backbone scale.} If depth-compounding drift is real, a deeper
backbone should drift more in late blocks. Re-running the main curve and CKA
with DINOv2 ViT-B/14 at pilot scale: clean accuracy \ViTBAccClean{} and the
same floor (\ViTBAccFloor{} at $s{=}5$), with the late-minus-early CKA gap at
\ViTBLateMinusEarly{} --- the gradient replicates and steepens
(\ViTBLate{} vs.\ \ViTBEarly{} on ViT-B against \LateDrift{} vs.\
\EarlyDrift{} on ViT-S). The mechanism is not a ViT-S artifact.

\paragraph{Prior vs.\ geometry: a discriminating test.} The natural reading
of the constant-class collapse is that the readout replays its training
prior. We tested this directly with two interventions on the readout's
training distribution, holding the backbone and test fold fixed.
Class-balanced reweighting is a no-op on our near-uniform fold (collapse
unchanged: frog, 99\%). So we made the prior genuinely non-uniform:
oversampling a single non-frog class (cat) five-fold, so the training
prior's dominant class is unambiguous. If the collapse replayed the prior,
the darkest fold should now predict ``cat.'' It does not: the sev-5 collapse
remains \emph{frog} at the same 99\% share and the same 0.02 entropy. The
constant class is not chosen by the readout's prior at all --- in the dark
embedding geometry, images of \emph{every} class land nearest the frog
prototype. The collapse target is a property of the degraded representation,
which also explains why it is invisible to any intervention on the readout's
training distribution.

\paragraph{Compute cost.} Frozen ViT-S/14 runs at 40.6\,ms per image on CPU
(batch 1; interleaved-round medians, Appendix~\ref{app:eff}); the LoRA arms
add 21--25\% latency for at
most 1.18M trainable parameters (5.3\% of the backbone). The full table,
including throughput, is \texttt{results\_pilot/efficiency.csv}; the point
for deployment is that the readout interventions of \S\ref{sec:mech} are free
relative to any backbone modification.

% ============================================================================
\section{Discussion}
\label{sec:discussion}

The measurements compose into one account. Photon loss degrades contrast and
SNR in the continuous image; the probe's predominantly low-frequency signal
weakens; drift compounds through depth rather than appearing at the input;
decision margins shrink until argmax crosses thresholds, which it does in
steps rather than smoothly; and once margins are gone, the readout falls back
on its training prior, which is why the darkest condition looks like a
constant-class classifier with a frog habit. Digitization, measured directly,
contributes nothing. The information a retrained readout needs is still in
the embeddings --- 22.5 points of it at severity 4 --- which is what makes
this a readout problem and not a representation-annihilation problem.

Three corollaries, each falsifiable with the released harness --- and one
already tested. First (\emph{tested, refuted, replaced}): we predicted the
floor's class identity would follow the readout's training prior; the
imbalanced-prior experiment of \S\ref{sec:robust} refutes this --- even with
a five-fold-oversampled non-frog prior, the collapse remains frog. The
revised claim: the collapse target is a geometry attractor determined by
which class prototype sits nearest the collapsed dark manifold, and it
should be predictable from embedding distances alone, before any readout is
trained. Second, severity-matched augmentation of the readout's training
fold should recover most of what the adapted probe recovers, without ever
seeing test-time corruptions at training time. Third, any intervention that
preserves late-block geometry --- whatever its form --- should beat any
intervention that operates purely in input space; this one is already
confirmed at pilot scale, where the adapted readout beats the best
enhancement by +15.8 points at $s{=}4$ ($p{=}0.007$) and +20.0 points at
$s{=}5$ ($p{<}10^{-3}$).

For practitioners, the guidance is short. Do not buy float-sensor capture for
a frozen DINOv2 pipeline; it purchases nothing. Do not bolt CLAHE onto the
front; it will cost you points on clean data. If you adapt, anchor the
objective --- and expect even that to be second-best to simply recalibrating
the readout on data that looks like deployment.

% ============================================================================
\section{Limitations}
\label{sec:limitations}

The honest ones. One backbone family (DINOv2) at two scales; one dataset
(CIFAR-10 upsampled from 32px, which flattens the frequency content that the
frequency controls probe); a synthetic, globally uniform photometric ladder
with no vignetting, auto-exposure dynamics, or real sensor MTF; and a
pilot-scale test fold, which makes several paired comparisons
under-powered --- we flag the ones that are ($s{=}3$ adaptation, $p{=}0.14$)
rather than narrating them as trends. The LoRA ablation runs one seed per
configuration; the negative effect is consistent across all six
configurations, but per-configuration variance is unmeasured until the
full-scale run. Real-darkness validation (ExDark \cite{wei2018deep}, SID
\cite{chen2018learning}) is the obvious next step and is config-flagged in
the harness.

% ============================================================================
\section{Conclusion}

Low-light failure in a frozen self-supervised ViT is a readout problem before
it is a fine-tuning problem. The damage happens in the continuous image, not
at the ADC; it compounds through depth; and the fixed readout answers it by
degenerating to a constant class --- quietly, in a way per-class accuracy
reports as a success story. Retraining the readout on degraded embeddings
recovers what enhancement, adapters, and sensor upgrades cannot. Every number
in this paper regenerates from one command against committed artifacts, and
the statistical machinery behind the claims ships with its own test suite.

% ============================================================================
{\small
\bibliographystyle{ieeenat_fullname}
\bibliography{references}
}

\appendix
\section{Reproducibility}
\label{app:repro}

Every experiment is a subcommand of \texttt{run\_experiments.py}
(\texttt{main\_curve}, \texttt{cka}, \texttt{frequency}, \texttt{mechanism},
\texttt{quantization}, \texttt{lora\_ablation}, \texttt{mitigation},
\texttt{efficiency}, \texttt{failure\_analysis}, \texttt{seed\_sensitivity},
\texttt{cross\_dataset}) driven by \texttt{configs/experiments.yaml}; the
config SHA-256, package versions, and dataset indices are written to
\texttt{results/manifest.jsonl} per run. The statistical machinery --- unbiased
CKA, Wilson and bootstrap intervals, permutation tests, Benjamini--Hochberg
correction, the two curve-shape nulls --- is covered by 22 unit tests
(\texttt{tests/run\_tests.py}), including a verification that the biased CKA
estimator inflates similarity at $n \ll d$ while the estimator used here does
not. Noise draws are derived per-image (hash of image bytes and severity), so
corruption is bit-reproducible across scripts, processes, and experiment
orders. The LoRA ablation, anchored-adapter arm, and STL-10 transfer run on a
single GPU via \texttt{run\_all\_colab.py} and fill
Tables~\ref{tab:lora}--\ref{tab:lora-anchored}.

\section{Efficiency Detail}
\label{app:eff}

\begin{table}[h]
\centering\small
\caption{Latency (CPU, batch 1, interleaved-round medians over 4 rounds of
25 iterations each).}
\label{tab:eff}
\InputIfFileExists{results/table_eff}{}{\textit{[Efficiency table pending.]}}
\end{table}

\end{document}
